\documentclass[11pt]{article}

\usepackage[margin=25mm]{geometry}

\usepackage{amsmath,amssymb,fontspec,unicode-math}

\IfFontExistsTF{DejaVu Serif}{\setmainfont{DejaVu Serif}}{\setmainfont{Latin Modern Roman}}

\IfFontExistsTF{Latin Modern Math}{\setmathfont{Latin Modern Math}}{\setmathfont{DejaVu Math TeX Gyre}}

\usepackage{xurl,hyperref}

\setlength{\parindent}{0pt}\setlength{\parskip}{6pt}

\begin{document}

\title{Revisiting Bi-Linear State Transitions in Recurrent Neural Networks}

\author{ELVIA preliminary AI analysis}\date{21 October 2026}\maketitle

Preliminary assessment: Minor.\par

Independent scientific review remains pending. Author adoption has not been requested.

The unrestricted-length parity probability counts a length-dependent sign orientation, while two examples do not identify an unspecified learning rule. Orthogonal canonical-form and positive-diagonal state-tracking explanations also require narrower scope. The main constructive representation results survive.

\section*{Review scope}

Claim-focused reading of the original paper and relevant author-released source, with direct mathematical and static-source checks. No author program, historical verifier, model training or published benchmark was executed.

\section*{Problem 1: The counted parity orientation depends on length}

Source: Proposition 3, p.6; Appendix A.3, pp.15–16 and final author remark

Assessment: supported; impact: minor

\subsection*{Source claim (paraphrase)}

The proposition claims arbitrary-length parity from two examples with probability \(1-2^{-H}\), counting either opposite-sign transition orientation.

\subsection*{Analysis}

A coordinate equals \(h_{0,i}(a_i^{[0]})^{T-k}(a_i^{[1]})^k\). Signs \((+,-)\) encode \((-1)^k\), but \((-,+)\) encode \((-1)^T(-1)^k\). For \(H=1,h_0=1\) and independent symmetric Rademacher entries, only \((+,-)\) succeeds for all lengths. The other outcomes have opposite-label hidden-state collisions: \((+,+)\) and \((-,-)\) collide on inputs zero/one, while \((-,+)\) collide on one/one-one. Thus success is \(1/4\), not claimed \(1/2\). With two coordinates \(A_0=\operatorname{diag}(1,-1),A_1=\operatorname{diag}(-1,1)\), training inputs zero-zero/zero-one permit either coordinate readout; only the first extrapolates to odd lengths. Existence therefore does not identify an arbitrary trained readout.

\subsection*{Finding and scope}

The probability and learning interpretation need qualification. The author’s final remark already gives the corrected sufficient-event probability; even-only experimental lengths do not test the unrestricted claim.

\subsection*{Proposed correction}

Credit the existing author remark: with independent symmetric nonzero entries and nonzero initialization, a selected \((+,-)\) coordinate exists with probability \(1-(3/4)^H\). Known matrices can identify it and two opposite labels orient it. This is a sufficient-event constructive rule, not an exact generic-optimizer probability; fixed length parity separately permits either orientation.

The final Appendix A.3 remark already states the appropriate sufficient-event probability. The proposition and unrestricted two-example learning interpretation remain broader.

\section*{Problem 2: Even orthogonal dimension does not remove reflection blocks}

Source: Section 2.3 and Appendix A.2, rotation-block explanation

Assessment: supported; impact: minor

\subsection*{Source claim (paraphrase)}

The explanation represents an arbitrary even-dimensional real orthogonal matrix entirely by planar rotation blocks.

\subsection*{Analysis}

Take \(A=\operatorname{diag}(1,-1)\). It is orthogonal in dimension \(2\) but determinant \(-1\); every rotation-block-only matrix has determinant \(1\), preserved by similarity. The general real form also includes one-dimensional \(+1,-1\) blocks. The paper’s existence construction explicitly chooses rotations, so it does not need the false universal characterization.

\subsection*{Finding and scope}

The canonical-form explanation is overbroad; the cyclic and finite-product representation construction remains valid.

\subsection*{Proposed correction}

Use the real orthogonal form with \(+1,-1\) eigenspaces and conjugate rotation blocks. For a single even-dimensional determinant-one matrix those real signs can be paired; a family additionally needs a compatible common basis.

A proposed correction is described; author adoption is not established.

\section*{Problem 3: Positive diagonality still permits absorbing-state tracking}

Source: Section 1 hierarchy and Section 2.1 finite-state simulation definition

Assessment: supported; impact: minor

\subsection*{Source claim (paraphrase)}

A parity obstruction should not be interpreted as absence of every nontrivial finite-state tracking task.

\subsection*{Analysis}

Let the state report whether any one has appeared. With \(h_0=(1,1)\), \(A_0=\operatorname{diag}(1/2,1/2)\), \(A_1=\operatorname{diag}(1/2,1/4)\), all entries are strictly positive contractions. After length \(T\) with \(k\) ones, \(h_T=2^{-T}(1,2^{-k})\). Scores \(2^{-T-k}\) for unseen-one state and \((3/4)2^{-T}\) for seen-one state classify correctly for \(k=0\) versus \(k\ge1\), including empty input. No additive term or length-aware decoder is used.

\subsection*{Finding and scope}

The broad verbal hierarchy needs narrowing to the forbidden cyclic/group tasks. This exact-arithmetic construction proves neither parity nor robust unlimited-length machine arithmetic nor trainability.

\subsection*{Proposed correction}

Name the specific excluded tasks and distinguish universal finite-state simulation from particular absorbing-state machines. No complete classification of positive-diagonal automata is supplied.

A proposed correction is described; author adoption is not established.

\section*{Limits}

Preliminary AI assessment; independent human scientific adjudication remains pending.\par

Analytic examples and release discrepancies do not remeasure the paper’s reported experiments.\par

\section*{Source identity}

Current official NeurIPS 2025 publisher PDF, freshly retrieved 21 October 2026

\url{https://proceedings.neurips.cc/paper_files/paper/2025/file/64d71dc33dbd217e3752c9ff1a41798d-Paper-Conference.pdf}

Reviewed PDF SHA-256: \nolinkurl{46f9079250d93617468af56159c232c65d0f010313b554531c27b3265727f2e3}

Fresh HTTP 200 download from the official publisher; actual retrieval timestamp and SHA-256 recorded. The content is byte-identical to the retained proceedings PDF. No file modification time is used as provenance.

\end{document}
